% =============================================================================
% Relatorio 2 (PRO3474, Grupo 4, parceria Riachuelo): estudo de diferenciacao
% (item B do roteiro). Conteudo: metodo, referencias comparadas (Tabela
% tab:r2-diferenciacao), mapa de posicionamento (Figura fig:r2-posicionamento),
% diferenciais DF1 a DF5, o que nao e diferencial e riscos.
% Nao repete precos (Secao sec:r2-valor) nem solucoes tecnicas a aproveitar
% (Secao sec:r2-aproveitamento). Posicoes do mapa: estimativa qualitativa da equipe.
% =============================================================================

\section{ESTUDO DE DIFERENCIAÇÃO}
\label{sec:r2-diferenciacao}

As funções da Seção~\ref{sec:r2-funcional} dizem o que o Tag\&Go precisa fazer, mas não em que ele se distingue do que já existe. O benchmarking competitivo do Relatório 1, que deixou a Riachuelo com 2,67, atrás da C\&A, da Zara e da Renner (Subseção~\ref{sub:r2-sintese}), respondeu só em parte a essa pergunta. Aquela comparação descrevia a rede tradicional da Riachuelo e considerava apenas varejistas de moda, sem incluir as soluções que já permitem ao cliente pagar e sair sem passar pelo caixa. Como o Tag\&Go será vendido a um varejista, a equipe precisava saber em que ele se distingue tanto dos concorrentes da Riachuelo quanto das empresas que vendem tecnologia de autoatendimento, e também da própria Riachuelo, que abriu em 2026 lojas conceito com self-checkout. A equipe fez o estudo de diferenciação para responder a essa pergunta, e o resultado alimentou as decisões de concepção da Seção~\ref{sec:r2-concepcao}.

\subsection{Método}
\label{sub:r2-dif-metodo}

A equipe comparou o Tag\&Go com dez referências em sete atributos: onde o cliente paga; onde a peça é liberada; se existe barreira física que só abre com um pagamento válido; se é preciso instalar um aplicativo; se há programa de fidelidade e uso de dados do cliente; se a etiqueta volta à origem para ser reutilizada; e em que estágio e região a solução está. Os atributos foram escolhidos a partir da casa da qualidade do Relatório 1, em que a liberação condicionada ao pagamento (R06) concentrou 33,1\% da importância técnica, o tempo de pagamento (R03), 24,4\%, e a compra sem funcionário (R04), 13,2\%. Os dois primeiros atributos traduzem R03 e R04, o terceiro traduz R06, e os demais descrevem o que o cliente precisa ter no celular, o que o varejista ganha em relacionamento e o que acontece com a etiqueta depois da venda.

As referências foram reunidas em três grupos. O primeiro é a própria Riachuelo, em dois recortes, a rede tradicional e as lojas conceito de 2026, porque a nota do Relatório 1 não descreve o formato novo. O segundo reúne os concorrentes brasileiros já usados no QFD: Renner, C\&A e Zara. O terceiro reúne soluções que liberam a peça ou a saída sem caixa em outros mercados: MishiPay, Amazon Just Walk Out com RFID, Paytag, rapitag e Sensormatic SuperTag 4. As informações vieram de páginas das empresas, de notícias e de estudos de caso publicados; os números divulgados pelos próprios fornecedores foram atribuídos a eles no texto e na tabela.

Dois recortes ficaram fora do estudo. Os preços das soluções entram na escala vertical da Seção~\ref{sec:r2-valor}, e as soluções técnicas que o Tag\&Go pode aproveitar dessas referências, inclusive as patentes, entram no estudo de aproveitamento técnico da Seção~\ref{sec:r2-aproveitamento}. O estudo tratou do que cada solução entrega ao cliente e ao varejista.

\subsection{Referências comparadas}
\label{sub:r2-dif-referencias}

Na rede tradicional da Riachuelo, o cliente paga no caixa e o funcionário retira a etiqueta rígida com um destacador magnético. A etiqueta é magnética comum, de modo que qualquer destacador abre a trava, e o tempo de pagamento foi estimado em 120\,s no Relatório 1 (estimativa da equipe). A rede usa RFID com a Mozaiko desde 2021, da expedição da fábrica de Natal às lojas, e, segundo o presidente da empresa, cada produto tem um chip \cite{tiinside2022riachuelo,neofeed2025riachuelo}; a rede também mantém o Ria Club, programa de fidelidade por cashback \cite{midway2026riaclub}. O app da marca passa de 10 milhões de downloads, mas não localiza a peça na loja nem mostra o estoque por loja \cite{googleplay2026riachuelo,appstore2026riachuelo}.

As lojas conceito mudaram parte desse quadro. A loja de Curitiba, inaugurada em 22/04/2026, tem self-checkout com etiquetas inteligentes em todos os pisos \cite{amanha2026riachuelo}; a da Oscar Freire, reaberta em 23/09/2026, combina provadores inteligentes, self-checkout e etiquetas inteligentes \cite{economicnews2026riachuelo}; e o pop-up de Pinheiros, de dezembro de 2025, opera com 100\% de self-checkout por RFID \cite{neofeed2025riachuelo}. A empresa planeja de 15 a 20 lojas no formato em 2026 e pretende levar o conceito a mais de 300 unidades \cite{mercadoconsumo2026riachuelo}. A tecnologia da etiqueta inteligente não foi divulgada, e por isso a equipe não sabe se a etiqueta das lojas conceito tem barreira autenticada ou se é aberta no terminal por um destacador. Nas duas hipóteses, o pagamento e a liberação continuam presos a um terminal fixo.

Entre os concorrentes brasileiros, a Renner tem RFID em 100\% das mercadorias e estoque atualizado a cada 60\,s \cite{infomoney2024renner}. Seus 742 terminais de autoatendimento por RFID estão em 213 lojas, 52\% da rede, e respondem por cerca de 38\% dos pagamentos onde existem \cite{consumidormoderno2025renner}. A leitura e a desativação das etiquetas acontecem em lote no terminal, e a proteção na saída é feita pelo RFID usado como EAS, sem barreira que dependa do pagamento. A C\&A seguiu caminho parecido no conceito Energia, com self-checkout por RFID \cite{brazileconomy2025cea}, e acrescentou o pagamento por biometria facial para clientes do C\&A Pay, que exige cadastro facial no app \cite{mobiletime2023cea}; o C\&A Pay devolve 10\% em cashback \cite{cea2026pay}. A Zara oferece o Store Mode, com estoque da loja, retirada em 30 minutos e localização da peça na planta \cite{inditex2021storemode}. Em 2023, a empresa atrasou a implantação de um sistema de alarme por RFID sem etiqueta rígida, porque os chips eram fáceis de arrancar e o alarme tocava em devoluções \cite{pymnts2023zara}. A Inditex também é o exemplo mais conhecido de recirculação: desde 2014, suas etiquetas rígidas com RFID e AM voltam a centros da Tyco para reuso \cite{rfidjournal2014inditex}.

As soluções do terceiro grupo resolvem a saída de formas distintas. Na MishiPay, o cliente lê o código da peça no celular, paga com Apple Pay ou Google Pay, e um portal RFID na saída confere os identificadores pagos numa lista branca; segundo o fornecedor, a solução está em mais de 2.000 lojas \cite{mishipay2026apparel}. Na Flying Tiger, a mesma empresa oferece a compra num web app aberto por QR, sem instalar aplicativo \cite{mishipay2026flyingtiger}. A Amazon adotou RFID no Just Walk Out para vender roupas, porque a visão computacional exige produtos em prateleira, e usa um portal na saída com cartão ou palma da mão \cite{amazon2024jwo}. Nenhuma das duas mantém uma trava física na peça: ambas dependem de detectar a saída de um item não pago.

As três últimas referências mantêm a barreira física. Na Paytag, empresa israelense fundada em 2022, o cliente lê a etiqueta com o celular, paga no app e leva as peças a uma estação que remove as etiquetas depois do pagamento \cite{calcalist2026paytag}. A Sensormatic SuperTag 4 integra um destacador motorizado ao self-checkout, confere se o item pago é o da etiqueta e, segundo o fornecedor, só destaca com uma compra válida \cite{sensormatic2026supertag}. A rapitag, de Munique, é a referência mais próxima do Tag\&Go: o cliente lê o produto e paga no app, e a própria etiqueta abre. A empresa fez um piloto de três meses em 2018, no PEP de Munique, com fones de ouvido protegidos por etiquetas BLE \cite{storesshops2019rapitag} e outro em 2024, com um varejista alemão, em seis lojas e mais de 3.500 etiquetas, que abrem também no caixa e no self-checkout \cite{eurocis2024rapitag}; em todos os casos, o cliente precisa do app do varejista. A Tabela~\ref{tab:r2-diferenciacao} reúne as referências e o Tag\&Go nos sete atributos.

\begin{landscape}
\begin{table}[htbp]
    \centering
    \caption{Comparação do Tag\&Go com soluções de mercado e com a própria Riachuelo}
    \label{tab:r2-diferenciacao}
    \scriptsize
    \setlength{\tabcolsep}{3pt}
    \renewcommand{\arraystretch}{1.25}
    \begin{tabularx}{\linewidth}{@{}
        >{\RaggedRight\arraybackslash}p{2.5cm}
        >{\RaggedRight\arraybackslash}X
        >{\RaggedRight\arraybackslash}X
        >{\RaggedRight\arraybackslash}X
        >{\RaggedRight\arraybackslash}p{2.4cm}
        >{\RaggedRight\arraybackslash}X
        >{\RaggedRight\arraybackslash}p{2.6cm}
        >{\RaggedRight\arraybackslash}X@{}}
        \toprule
        \textbf{Referência} & \textbf{Onde paga} & \textbf{Onde libera} & \textbf{Barreira física autenticada} & \textbf{Instalar app} & \textbf{Fidelidade e dados} & \textbf{Etiqueta volta à origem} & \textbf{Estágio} \\
        \midrule
        Riachuelo, rede tradicional & caixa & caixa, com destacador magnético & não (etiqueta magnética comum) & não & Ria Club, por cashback & não documentado & operação; tempo de pagamento estimado em 120\,s \\
        Riachuelo, lojas conceito & terminal fixo de self-checkout & no terminal (tecnologia da etiqueta inteligente não divulgada) & não divulgado & não & Ria Club & não divulgado & Curitiba e Oscar Freire (2026); pop-up de Pinheiros (2025) \\
        Renner & terminal RFID fixo (cerca de 38\% dos pagamentos onde existe) & leitura e desativação em lote no terminal & não (EAS por RFID) & não & cartão Realize; CRM com IA & não documentado & 742 terminais em 213 lojas \\
        C\&A (Energia) & terminal RFID; pagamento facial no C\&A Pay & no terminal & não & cadastro facial no app & C\&A Pay, cashback de 10\% & não documentado & lojas conceito \\
        Zara & self-checkout e app (Store Mode) & self-checkout & não; sistema sem etiqueta rígida atrasou em 2023 & app para o Store Mode & app & recirculação com a Tyco (2014) & global \\
        MishiPay & celular (leitura e pagamento) & sem etiqueta rígida; portal RFID com lista branca & não & app ou web app aberto por QR & app do varejista & não se aplica & mais de 2.000 lojas, segundo o fornecedor \\
        Amazon Just Walk Out com RFID & portão com cartão ou palma da mão & portal RFID & não & não & Amazon & não se aplica & arenas e lojas de estádio \\
        Paytag & celular & estação de self-checkout que remove as etiquetas & sim, na estação & app & app & não documentado & Israel \\
        rapitag & celular & a própria etiqueta, pelo celular & sim & app do varejista & leasing e dados cogitados & não documentado & pilotos em 2018 e 2024 (Alemanha) \\
        Sensormatic SuperTag 4 & self-checkout & destacador motorizado no terminal & sim (só destaca com compra válida, segundo o fornecedor) & não & não & recirculação Sensormatic & comercial \\
        \midrule
        \rowcolor{r2laranjaclaro}
        Tag\&Go & celular, em qualquer ponto da loja & a própria etiqueta, pelo celular; Liberador de Balcão nas exceções & sim (HMAC verificado na etiqueta; carretel não ferromagnético) & não na Compra Rápida; app na Jornada Completa & Ria Club, recomendação restrita ao estoque da loja, métricas por categoria & sim, com recarga e cadastro na origem & conceito (Relatório 2) \\
        \bottomrule
    \end{tabularx}
    \source{Elaborado pelos autores (2026) com base em \citeonline{abvtex2026riachuelo,amanha2026riachuelo,economicnews2026riachuelo,neofeed2025riachuelo,consumidormoderno2025renner,infomoney2024renner,brazileconomy2025cea,mobiletime2023cea,inditex2021storemode,pymnts2023zara,rfidjournal2014inditex,mishipay2026apparel,mishipay2026flyingtiger,amazon2024jwo,calcalist2026paytag,storesshops2019rapitag,eurocis2024rapitag,sensormatic2026supertag}. O tempo de 120\,s da rede tradicional é estimativa da equipe, feita no Relatório 1.}
\end{table}
\end{landscape}

Lida por colunas, a tabela permitiu à equipe localizar o espaço livre. Todas as referências brasileiras, inclusive as lojas conceito da Riachuelo, concentram pagamento e liberação num terminal fixo, e nenhuma delas documenta uma barreira que dependa do pagamento. As soluções que dispensam o terminal, MishiPay e Amazon, abrem mão da trava na peça e confiam no portal da saída, o que em moda esbarra no mesmo problema que atrasou a Zara: sem etiqueta rígida, o controle depende de detectar o furto depois que ele ocorreu. As que mantêm a trava, Paytag e SuperTag 4, voltam a exigir uma estação. Só a rapitag combina liberação na peça e pagamento no celular, e ela exige o app do varejista. Na coluna da etiqueta, por fim, a equipe constatou que a recirculação de etiquetas rígidas já é prática de mercado (Zara com a Tyco e a própria Sensormatic), mas sempre com etiquetas passivas.

\subsection{Mapa de posicionamento}
\label{sub:r2-dif-mapa}

Para sintetizar a comparação, a equipe posicionou as referências num mapa com dois eixos, derivados dos atributos de maior peso no QFD. O eixo horizontal mede a independência de ponto fixo para concluir a compra, de 0 (caixa com funcionário) a 10 (nenhum ponto fixo), e representa R03 e R04, porque cada deslocamento até um terminal ou estação acrescenta tempo e uma chance de fila. O eixo vertical mede a barreira física na peça que só abre com pagamento válido, de 0 (nenhuma) a 10 (barreira autenticada na própria peça), e representa R06. Esse segundo eixo pesa na decisão do varejista: no estudo do ECR de 2026, com 39 varejistas, 35 deles supermercados, as lojas com self-checkout perderam 33\% mais que as demais \cite{ecr2026selfcheckout}, e cerca de 80\% dos decisores ouvidos pela Zebra veem o furto como grande problema no self-checkout \cite{zebra2023shopper}.

As posições foram atribuídas pela equipe e são uma estimativa qualitativa, sem medição. No eixo horizontal, o caixa recebeu 1; os terminais fixos de autoatendimento e as estações, de 2,5 a 4, conforme a parte da jornada que já acontece no celular; as soluções em que o cliente conclui a compra sem passar por terminal ou estação, de 7,5 a 8 (MishiPay, Amazon e rapitag); e o Tag\&Go, 9, porque ainda depende de um ponto fixo, o Liberador de Balcão, nas exceções. No eixo vertical, as soluções sem etiqueta rígida receberam 2 ou 2,5; as etiquetas rígidas comuns, abertas por qualquer destacador, de 3 a 4,5; as barreiras autenticadas numa estação, 7,5 ou 8; e as barreiras na própria peça, 8 (rapitag) e 9 (Tag\&Go). O Tag\&Go não recebeu 10 porque nem a resistência à abertura sem pagamento (R06) nem a imunidade ao destacador (E11) foram ensaiadas; as duas serão verificadas no protótipo. A Figura~\ref{fig:r2-posicionamento} mostra o resultado.

\begin{figure}[htbp]
    \centering
    \caption{Mapa de posicionamento}
    \label{fig:r2-posicionamento}
    \begin{tikzpicture}
        \begin{axis}[
            r2eixo,
            width=0.95\textwidth,
            height=0.78\textwidth,
            xmin=0, xmax=10, ymin=0, ymax=10,
            xtick={0,1,2,3,4,5,6,7,8,9,10},
            ytick={0,1,2,3,4,5,6,7,8,9,10},
            xlabel={Independência de ponto fixo para concluir a compra\\(0 = caixa com funcionário; 10 = nenhum ponto fixo)},
            ylabel={Barreira física na peça que só abre com pagamento válido\\(0 = nenhuma; 10 = autenticada na própria peça)},
            xlabel style={align=center},
            ylabel style={align=center},
            clip=false,
        ]
            \fill[r2laranjaclaro] (axis cs:5,5) rectangle (axis cs:10,10);
            \draw[dashed, r2cinza] (axis cs:5,0) -- (axis cs:5,10);
            \draw[dashed, r2cinza] (axis cs:0,5) -- (axis cs:10,5);
            \node[r2rotulo, anchor=north west] at (axis cs:5.1,9.9) {sem ponto fixo e com\\barreira autenticada};
            \node[r2rotulo, anchor=north west] at (axis cs:0.1,9.9) {barreira autenticada\\em ponto fixo};
            \node[r2rotulo, anchor=south west] at (axis cs:0.1,0.1) {caixa ou terminal fixo,\\sem barreira autenticada};
            \node[r2rotulo, anchor=south east] at (axis cs:9.9,0.1) {sem ponto fixo,\\sem barreira na peça};
            \addplot[only marks, mark=*, mark size=2.5pt, color=r2azulmedio]
                coordinates {(1,4) (3,3) (3,3.5) (3.5,4) (2.5,8) (4,7.5) (7.5,2) (8,2.5) (8,8)};
            \addplot[only marks, mark=o, mark size=2.5pt, thick, color=r2azulmedio]
                coordinates {(2.5,4.5)};
            \addplot[only marks, mark=*, mark size=3.5pt, color=r2laranja]
                coordinates {(9,9)};
            \node[font=\scriptsize, align=center, text=r2azul, anchor=north] at (axis cs:1,3.8) {Riachuelo\\(rede tradicional)};
            \node[font=\scriptsize, align=center, text=r2azul, anchor=south] at (axis cs:2.5,4.7) {Riachuelo\\(lojas conceito)};
            \node[font=\scriptsize, align=center, text=r2azul, anchor=west] at (axis cs:3.15,3) {Renner};
            \node[font=\scriptsize, align=center, text=r2azul, anchor=west] at (axis cs:3.15,3.5) {C\&A Energia};
            \node[font=\scriptsize, align=center, text=r2azul, anchor=west] at (axis cs:3.65,4) {Zara};
            \node[font=\scriptsize, align=center, text=r2azul, anchor=north] at (axis cs:2.5,7.8) {Sensormatic\\SuperTag 4};
            \node[font=\scriptsize, align=center, text=r2azul, anchor=west] at (axis cs:4.15,7.5) {Paytag};
            \node[font=\scriptsize, align=center, text=r2azul, anchor=north] at (axis cs:7.5,1.8) {MishiPay};
            \node[font=\scriptsize, align=center, text=r2azul, anchor=west] at (axis cs:8.15,2.5) {Amazon JWO\\com RFID};
            \node[font=\scriptsize, align=center, text=r2azul, anchor=north] at (axis cs:8,7.8) {rapitag\\(exige app)};
            \node[font=\scriptsize, align=center, text=r2laranja, anchor=south] at (axis cs:9,9.2) {Tag\&Go};
        \end{axis}
    \end{tikzpicture}
    \source{Elaborado pelos autores (2026), com posições estimadas pela equipe a partir das fontes da Tabela~\ref{tab:r2-diferenciacao}. Marcador vazado: tecnologia da etiqueta não divulgada.}
\end{figure}

O quadrante superior direito, que reúne liberação sem ponto fixo e barreira autenticada, tem apenas duas soluções: a rapitag, que exige o app do varejista, e o Tag\&Go. O quadrante inferior esquerdo concentra o varejo de moda brasileiro, inclusive a Renner, líder no benchmarking do Relatório 1, e daí a equipe concluiu que o self-checkout em terminal fixo virou prática do segmento e já não distingue quem o adota. As lojas conceito da Riachuelo aparecem com marcador vazado porque sua posição vertical é incerta: se a etiqueta inteligente tiver barreira autenticada, o ponto sobe em direção à SuperTag 4, mas continua à esquerda, preso ao terminal. A distância entre o Tag\&Go e a rede tradicional da Riachuelo, nos dois eixos, mede o salto que o sistema propõe à empresa.

\subsection{Diferenciais do Tag\&Go}
\label{sub:r2-dif-diferenciais}

A partir da tabela e do mapa, a equipe identificou cinco diferenciais, codificados de DF1 a DF5 para não se confundirem com os desejos D01 a D07 do Relatório 1. Cada diferencial foi definido pelo contraste com uma referência concreta, e não pela descrição do próprio produto.

O primeiro diferencial (DF1) é o pagamento em qualquer ponto da loja com liberação na própria peça, sem terminal fixo no fluxo principal. Frente às lojas conceito da Riachuelo, à Renner e à C\&A, o Tag\&Go elimina o ponto de pagamento; frente à Paytag e à SuperTag 4, elimina a estação de remoção. O único ponto fixo que resta é o Liberador de Balcão, usado nas exceções, como celular sem bateria ou etiqueta sem energia, e é ele que garante a alternativa humana exigida por R07.

O segundo diferencial (DF2) é a autorização verificada na própria etiqueta. A etiqueta gera um desafio a cada liberação, um nonce, e só aceita um token calculado para aquele desafio, aquela etiqueta e aquela transação; o celular apenas retransmite o token e não guarda segredo, de modo que um aparelho adulterado não consegue forjá-lo. Com o carretel não ferromagnético, a etiqueta também fica imune ao destacador magnético comum, vendido livremente por cerca de R\$~300 \cite{canalautomacao2026desacoplador}, que abre tanto a etiqueta da rede tradicional quanto qualquer etiqueta rígida comercial. O protocolo completo está na Seção~\ref{sec:r2-arquitetura}. Frente à rapitag, única outra referência com barreira na peça, a equipe não comparou o esquema de autorização, porque não o verificou; o DF2 foi afirmado, portanto, frente às demais referências.

O terceiro diferencial (DF3) é a Compra Rápida sem instalar aplicativo nas duas plataformas, pelo App Clip Tag\&Go no iPhone e pela Página de Compra Rápida no Android, combinada com a barreira física na peça. Nas duas plataformas, o cliente paga pela carteira do celular, com Apple Pay ou Google Pay e confirmação por Face ID, Touch ID ou senha na folha do sistema, ou por Pix copia e cola, sem cadastro prévio. A rapitag, a Paytag e a C\&A exigem app, e a Zara exige o app para o Store Mode. O web app da MishiPay na Flying Tiger é o precedente mais próximo de compra sem app \cite{mishipay2026flyingtiger}, só que sem etiqueta rígida. O diferencial tem duas dependências que a equipe declara desde já: no iPhone, o uso de Core Bluetooth dentro do App Clip foi deduzido da documentação da Apple, que não o proíbe, e será o primeiro teste de go/no-go do protótipo \cite{apple2026appclipfuncoes}; e o App Clip exige que o app Riachuelo publique a mesma função \cite{apple2026appclipcriar}. No Android, a página depende do Web Bluetooth, disponível no Chrome para Android e no Samsung Internet \cite{chrome2026webbluetooth}.

O quarto diferencial (DF4) é o ciclo fechado de uma etiqueta ativa, com cadastro na origem, no ramo fábrica (peças próprias, em Natal) ou no ramo CD (peças de terceiros), e recarga no CD, de modo que a loja não etiqueta nem cadastra peças. A recirculação de etiquetas passivas já existe: a Sensormatic recolhe etiquetas e pinos no ponto de venda e os devolve às fábricas \cite{sensormatic2024recirculacao}, a Pact relata reuso de até 80\% \cite{pact2026etiquetas}, e a Inditex recircula as suas com a Tyco \cite{rfidjournal2014inditex}. O que muda no Tag\&Go é fazer isso com uma etiqueta que tem bateria e chave criptográfica, o que acrescenta ao ciclo a inspeção, a recarga e a associação da etiqueta ao EPC da peça, descritas na Subseção~\ref{sub:r2-ciclo}. O trabalho de aplicar a etiqueta não desaparece; ele sai da loja e passa à origem. A meta de retorno de 99\% por ciclo (S01) fica acima dos 80\% a 83\% observados no mercado (Subseção~\ref{sub:r2-especificacoes}) e é tratada como meta, com os mecanismos e as consequências econômicas discutidos na Seção~\ref{sec:r2-valor}.

O quinto diferencial (DF5) é a Jornada Completa ligada ao Ria Club, no módulo Tag\&Go do app Riachuelo: um perfil com tamanhos e preferências, recomendação só do que existe na loja no tamanho do perfil, compra on-line do tamanho ausente, com retirada na loja sem custo ou entrega sem frete acima de um valor mínimo a definir com a Riachuelo, e métricas de sustentabilidade apresentadas como estimativa por categoria, com o cashback liberado depois que a etiqueta é devolvida ao Coletor de Etiquetas. Como o app Riachuelo não localiza a peça na loja nem mostra o estoque por loja, a Jornada Completa aproveita a base RFID que a rede já tem para preencher essa lacuna. A localização da peça entra como função secundária, no nível N1, o setor, obtido do estoque RFID e de um planograma digital, sem depender da etiqueta (Subseção~\ref{sub:r2-localizar}). Trata-se de um diferencial de experiência e de dados, sem relação com o hardware, e apoiado em hipóteses: o survey do Relatório 1 não mediu interesse por fidelidade, recomendação ou sustentabilidade.

A equipe também registrou o que não é diferencial, para não vender como novidade o que o mercado já oferece. O self-checkout em si já existe na Renner, na C\&A e nas lojas conceito da Riachuelo; o RFID, segundo a própria empresa, já está em todos os produtos da Riachuelo \cite{neofeed2025riachuelo}; o pagamento por carteira digital é oferecido pela MishiPay e por qualquer loja que aceite Apple Pay ou Google Pay; a localização da peça na planta da loja é feita pela Zara \cite{inditex2021storemode}; e o cashback já está no Ria Club e no C\&A Pay \cite{midway2026riaclub,cea2026pay}. Esses elementos entram no Tag\&Go como parte do núcleo comum ou da Jornada Completa, e a diferenciação vem da forma como o sistema os combina com a liberação na peça.

\subsection{Riscos de diferenciação}
\label{sub:r2-dif-riscos}

A posição do Tag\&Go no mapa pode ser contestada por três caminhos. O primeiro é a entrada no Brasil das duas empresas mais próximas. A Paytag já tem uso comercial em Israel e afirma ter patente registrada globalmente \cite{calcalist2026paytag}, e a rapitag já fez um piloto com um varejista na Alemanha e divide com o Tag\&Go o quadrante superior direito do mapa. Contra a rapitag, o DF3 é o diferencial mais defensável, porque a compra sem app depende de recursos de plataforma que a equipe já mapeou; contra a Paytag, pesa o DF1. O segundo caminho é a própria Riachuelo: como a tecnologia das lojas conceito não foi divulgada, ela pode já cobrir parte do DF1, e a equipe precisa confirmar com a empresa o que a etiqueta inteligente faz antes de apresentar o Tag\&Go como substituto. O terceiro caminho é a propriedade intelectual, já que há patentes de etiquetas que se destacam depois de uma compra verificada no celular; a análise das famílias de patentes está na Seção~\ref{sec:r2-aproveitamento}, e a busca no INPI ficou para o Relatório 3.

Os diferenciais também têm durabilidades diferentes. O DF5 é o mais fácil de copiar, porque depende de software e de dados que qualquer varejista com app e RFID pode reunir, como já faz a Renner. O DF2 e o DF4 são os mais difíceis, porque exigem ao mesmo tempo a etiqueta ativa, o protocolo de autorização e a logística reversa com cadastro na origem. O DF1 e o DF3 ficam no meio: são visíveis para o cliente, e por isso mais fáceis de imitar, mas só funcionam com a barreira na peça; o DF3 depende ainda de a Riachuelo manter no app a função equivalente ao App Clip.

Por fim, a equipe ligou os diferenciais às notas do Relatório 1. Na casa da qualidade, a Riachuelo recebeu nota 2 na redução da espera no caixa (DE01) e no pagamento rápido sem perder o atendimento (DE05), as duas vozes que somaram 50,5\% do peso relativo. O DF1 e o DF3 atacam DE01, ao retirar o caixa e o terminal do fluxo principal; o Liberador de Balcão e o Ponto de Apoio preservam o atendimento pedido em DE05. Com esses diferenciais, o Tag\&Go mira nota 5 nas duas vozes, o que em DE05 fica acima do plano de qualidade de 4 fixado no Relatório 1. As metas quantitativas correspondentes, de tempo de pagamento, compra sem funcionário e clareza do fluxo, estão nas especificações-meta da Tabela~\ref{tab:r2-especificacoes} e serão verificadas no protótipo.
